% ============================================================================
% LANGENTWURF LE-01b: Sich vorstellen + Verb ser
% Spanisch Klasse 7 - Sequenz 1: Empezamos
% ============================================================================
% Tier: 1 (vollständig)
% Doppelstunde: 3-4 von 8
% Fachfarbe: #c41e3a (Karminrot)
% ============================================================================

\documentclass[11pt,a4paper]{article}

\usepackage[utf8]{inputenc}
\usepackage[T1]{fontenc}
\usepackage[ngerman]{babel}
\usepackage{geometry}
\usepackage{fancyhdr}
\usepackage{lastpage}
\usepackage{xcolor}
\usepackage{tcolorbox}
\usepackage{enumitem}
\usepackage{tabularx}
\usepackage{longtable}
\usepackage{array}
\usepackage{booktabs}
\usepackage{amssymb}

% ============================================================================
% KONFIGURATION
% ============================================================================

\newcommand{\fach}{Spanisch}
\newcommand{\fachfarbe}{c41e3a}
\newcommand{\jahrgangsstufe}{7}
\newcommand{\schulart}{Gesamtschule}
\newcommand{\stundenthema}{Sich vorstellen + Verb ser}
\newcommand{\reihenthema}{Empezamos}
\newcommand{\stundennummer}{1b}
\newcommand{\reihenstunden}{4}
\newcommand{\gerniveau}{A1}
\newcommand{\lehrwerk}{Qué pasa 1}
\newcommand{\lehrwerkseite}{S. 10-15}
\newcommand{\lerngruppengroesse}{24}

% ============================================================================
% LAYOUT
% ============================================================================
\geometry{left=2.5cm, right=2.5cm, top=2.5cm, bottom=2.5cm}
\definecolor{fach}{HTML}{\fachfarbe}
\definecolor{gray}{HTML}{666666}
\definecolor{lightgray}{HTML}{f5f5f5}
\definecolor{successgreen}{HTML}{2a7a4b}
\definecolor{warningorange}{HTML}{b35c00}
\definecolor{basis}{HTML}{2a7a4b}
\definecolor{standard}{HTML}{2c5aa0}
\definecolor{erweiterung}{HTML}{7b1fa2}

\tcbuselibrary{skins,breakable}

\newtcolorbox{infobox}[1][]{
    colback=lightgray,
    colframe=fach,
    fonttitle=\bfseries,
    title=#1,
    breakable
}

\newtcolorbox{scorebox}{
    colback=white,
    colframe=fach,
    fonttitle=\bfseries,
    title=Didaktik-Score,
    breakable
}

\pagestyle{fancy}
\fancyhf{}
\fancyhead[L]{\small\textcolor{gray}{\fach{} -- Jg. \jahrgangsstufe{} -- \stundenthema}}
\fancyhead[R]{\small\textcolor{gray}{LE-\stundennummer{} | Tier 1}}
\fancyfoot[C]{\small\thepage{} / \pageref{LastPage}}
\renewcommand{\headrulewidth}{0.4pt}

% Differenzierungsmarker
\newcommand{\B}{\textcolor{basis}{\textbf{[B]}}}
\newcommand{\Std}{\textcolor{standard}{\textbf{[S]}}}
\newcommand{\E}{\textcolor{erweiterung}{\textbf{[E]}}}

% ============================================================================
% DOKUMENT
% ============================================================================
\begin{document}

% ============================================================================
% DECKBLATT
% ============================================================================
\begin{titlepage}
    \centering
    \vspace*{2cm}
    {\large\textcolor{gray}{\schulart}}\\[2cm]
    {\Huge\textcolor{fach}{\textbf{Langentwurf}}}\\[0.5cm]
    {\large LE-\stundennummer{} | Tier 1 (vollständig)}\\[2cm]
    {\LARGE\textbf{\stundenthema}}\\[0.5cm]
    {\large Sequenz: \reihenthema}\\[2cm]
    \begin{tabular}{rl}
        \textbf{Fach:} & \fach{} (\gerniveau) \\[0.3cm]
        \textbf{Klasse:} & \jahrgangsstufe \\[0.3cm]
        \textbf{Lehrwerk:} & \lehrwerk, \lehrwerkseite \\[0.3cm]
        \textbf{Doppelstunde:} & \stundennummer{} (Std. 3-4 von 8) \\[0.3cm]
    \end{tabular}
    \vfill
\end{titlepage}

\tableofcontents
\newpage

% ============================================================================
% 1. BEDINGUNGSANALYSE
% ============================================================================
\section{Bedingungsanalyse}

\subsection{Lerngruppe}
\begin{tabular}{ll}
    Kursgröße: & \lerngruppengroesse{} SuS \\
    Jahrgangsstufe: & \jahrgangsstufe \\
    Sprachniveau: & \gerniveau{} (GER) -- Anfänger \\
    Fremdsprache: & 2. Fremdsprache (nach Englisch) \\
\end{tabular}

\subsection{Vorwissen aus 01a}
\begin{itemize}
    \item Begrüßungen: ¡Hola!, ¡Buenos días!, ¡Buenas tardes!
    \item Verabschiedungen: ¡Adiós!, ¡Hasta luego!
    \item Erste Strukturen: ¿Qué tal? -- Bien, gracias.
    \item Buchaufbau und Orientierung
\end{itemize}

\subsection{Differenzierung (intern)}
\begin{tabular}{|l|p{3.5cm}|p{3.5cm}|p{3.5cm}|}
\hline
& \B{} \textbf{Basis} & \Std{} \textbf{Standard} & \E{} \textbf{Erweiterung} \\
\hline
\textbf{Ziel} & BR: Rezeption & MR: Produktion mit Hilfe & GY: Freie Produktion \\
\hline
\textbf{ser} & Lückentext mit Vorgabe & Konjugation selbst & + Sätze bilden \\
\hline
\textbf{Dialog} & Nachsprechen & Mit Stichworten & Frei variieren \\
\hline
\end{tabular}

\vspace{0.5em}
\textbf{WICHTIG:} Diese Marker sind NUR für die Lehrkraft!

% ============================================================================
% 2. SACHANALYSE
% ============================================================================
\section{Sachanalyse}

\subsection{Sprachliche Strukturen}

\textbf{Verb \textit{ser} (Präsens):}
\begin{center}
\begin{tabular}{|l|l|l|}
\hline
\textbf{Person} & \textbf{Pronomen} & \textbf{ser} \\
\hline
1. Sg. & yo & soy \\
2. Sg. & tú & eres \\
3. Sg. & él/ella/usted & es \\
1. Pl. & nosotros/-as & somos \\
2. Pl. & vosotros/-as & sois \\
3. Pl. & ellos/ellas/ustedes & son \\
\hline
\end{tabular}
\end{center}

\textbf{Redemittel:}
\begin{itemize}
    \item \textbf{Name:} Me llamo... / Soy... / ¿Cómo te llamas?
    \item \textbf{Herkunft:} Soy de... / ¿De dónde eres?
    \item \textbf{Nationalität:} Soy alemán/alemana / español/española
\end{itemize}

\subsection{Kontrastive Betrachtung}

\begin{tabular}{|l|p{5cm}|p{5cm}|}
\hline
& \textbf{Deutsch} & \textbf{Spanisch} \\
\hline
Subjektpronomen & Obligatorisch & Optional (Verbendung zeigt Person) \\
\hline
Höflichkeitsform & Sie (3. Pl.) & usted/ustedes (3. Sg./Pl.) \\
\hline
Kopulaverb & nur \textit{sein} & \textit{ser} vs. \textit{estar} \\
\hline
\end{tabular}

\textbf{Typische Interferenz:} *``Yo soy soy Maria'' (Dopplung durch dt. Struktur)

\subsection{Kulturelle Aspekte}
\begin{itemize}
    \item Duzen (tú) ist unter Gleichaltrigen in Spanien üblich
    \item Siezen (usted) bei Respektspersonen, im Geschäftsleben
    \item In Lateinamerika: vosotros kaum verwendet, ustedes für 2. Pl.
\end{itemize}

% ============================================================================
% 3. DIDAKTISCHE ANALYSE
% ============================================================================
\section{Didaktische Analyse}

\subsection{Stellung in der Sequenz}
\begin{tabular}{|c|l|l|}
\hline
\textbf{Block} & \textbf{Thema} & \textbf{Status} \\
\hline
1a & Rally + Begrüßung & abgeschlossen \\
\textbf{1b} & \textbf{Sich vorstellen + ser} & \textbf{diese Stunde} \\
1c & Artikel/Plural + Alphabet & folgt \\
1d & ¿Qué tal? + Sequenztest & folgt \\
\hline
\end{tabular}

\subsection{Kompetenzziele (Rahmenplan MV)}
\begin{itemize}
    \item \textbf{Kommunikativ:} Sprechen (an Gesprächen teilnehmen), Hörverstehen
    \item \textbf{Sprachlich:} Grammatik (Verbkonjugation), Wortschatz (Vorstellung)
    \item \textbf{Interkulturell:} Duzen/Siezen in hispanischen Kulturen
    \item \textbf{Methodisch:} Dialogisches Lernen, Partnerarbeit
\end{itemize}

\subsection{Legitimation}
\textbf{Rahmenplan Spanisch MV (Kl. 7-10):}
\begin{itemize}
    \item Kompetenzbereich: Kommunikative Kompetenz -- Sprechen
    \item Standard: ``Die SuS können sich und andere vorstellen und einfache Fragen zur Person stellen und beantworten.''
\end{itemize}

% ============================================================================
% 4. STUNDENZIELE
% ============================================================================
\section{Stundenziele}

\subsection{Grobziel}
Die SuS erweitern ihre \textbf{kommunikative Kompetenz im Bereich Sprechen}, indem sie sich vorstellen, nach Namen und Herkunft fragen und das Verb \textit{ser} korrekt anwenden.

\subsection{Feinziele (operationalisiert)}
\begin{tabular}{|c|p{7.5cm}|c|c|}
\hline
\textbf{Nr.} & \textbf{Die SuS können...} & \textbf{AFB} & \textbf{Diff.} \\
\hline
FZ1 & die Formen von \textit{ser} (yo, tú, él/ella) \textbf{nennen} und korrekt aussprechen & I & alle \\
\hline
FZ2 & einen kurzen Vorstellungsdialog mit vorgegebenen Redemitteln \textbf{führen} & II & alle \\
\hline
FZ3 & die passende Form von \textit{ser} in Lückensätzen \textbf{einsetzen} und \textbf{begründen} & II & \Std{}/\E{} \\
\hline
FZ4 & einen eigenen Vorstellungsdialog \textbf{entwickeln} und kulturelle Unterschiede \textbf{erklären} & III & \E{} \\
\hline
\end{tabular}

\vspace{1em}
\textbf{AFB-Verteilung:} AFB I: 25\% | AFB II: 50\% | AFB III: 25\%

% ============================================================================
% 5. METHODISCHE ANALYSE
% ============================================================================
\section{Methodische Analyse}

\subsection{Einstieg}
\begin{itemize}
    \item \textbf{Methode:} Kurzvideo/Audio einer Vorstellung (authentisch)
    \item \textbf{Begründung:} Aktivierung, Hörverstehen, Neugier wecken
    \item \textbf{Alternative:} Lehrkraft stellt sich ``fremd'' vor
\end{itemize}

\subsection{Erarbeitung}
\begin{itemize}
    \item \textbf{Methode:} Entdeckendes Lernen -- ser-Formen aus Beispielen ableiten
    \item \textbf{Begründung:} Aktive Konstruktion statt passiver Rezeption
    \item \textbf{Scaffolding:} Farbcodierung der Endungen, Tabelle schrittweise füllen
\end{itemize}

\subsection{Übungsphasen}
\begin{itemize}
    \item \textbf{Methode:} Chorisches Sprechen $\rightarrow$ Partnerarbeit $\rightarrow$ Präsentation
    \item \textbf{Progression:} Gelenkt $\rightarrow$ Halbfrei $\rightarrow$ Frei
    \item \textbf{Differenzierung:} Rollenkarten mit/ohne Hilfen
\end{itemize}

\subsection{Medien}
\begin{itemize}
    \item Tafel/Whiteboard: Konjugationstabelle
    \item Lehrwerk: \lehrwerk, \lehrwerkseite
    \item Arbeitsblatt: AB-01b.pdf
    \item Lernpfad: LP-01b.html (Tablets, optional)
    \item Audio: Track 5-6 (CD1)
\end{itemize}

% ============================================================================
% 6. VERLAUFSPLANUNG
% ============================================================================
\section{Verlaufsplanung}

\textbf{1. Stunde (45 Min.)}

\begin{longtable}{|p{1cm}|p{1.5cm}|p{4cm}|p{3cm}|p{1.5cm}|p{2cm}|}
\hline
\textbf{Zeit} & \textbf{Phase} & \textbf{Lehrerhandeln} & \textbf{Schülerhandeln} & \textbf{SF} & \textbf{Material} \\
\hline
\endhead

0--5 & Einstieg & Begrüßung (span.), Audio abspielen: ``Me llamo Carlos. Soy de Madrid.'' & Hören, Vermutungen & Plenum & Audio \\
\hline

5--10 & Aktivier. & ``Was habt ihr verstanden?'' Redemittel sammeln & Nennen: Name, Ort & Plenum & Tafel \\
\hline

10--20 & Input & Verb \textit{ser} einführen: Beispiele an Tafel, Formen ableiten lassen & Muster erkennen, Tabelle ergänzen & Plenum & Tafel, AB \\
\hline

20--30 & Übung 1 & Chorisches Sprechen anleiten: ``Yo soy..., tú eres...'' & Nachsprechen, Aussprache üben & Plenum & -- \\
\hline

30--40 & Übung 2 & PA: Minidialoge mit Redemittel-Karten & Dialoge üben (Name, Herkunft) & PA & Karten \\
\hline

40--45 & Sicherung & 2-3 Paare präsentieren, Feedback & Präsentieren, zuhören & Plenum & -- \\
\hline
\end{longtable}

\vspace{1em}

\textbf{2. Stunde (45 Min.)}

\begin{longtable}{|p{1cm}|p{1.5cm}|p{4cm}|p{3cm}|p{1.5cm}|p{2cm}|}
\hline
\textbf{Zeit} & \textbf{Phase} & \textbf{Lehrerhandeln} & \textbf{Schülerhandeln} & \textbf{SF} & \textbf{Material} \\
\hline
\endhead

0--5 & Aktivier. & Schnelles Quiz: ``Wie heißt \textit{ich bin}?'' & Antworten: soy & Plenum & -- \\
\hline

5--15 & Vertiefung & Lückentext AB, ser einsetzen & AB bearbeiten & EA & AB-01b \\
\hline

15--25 & Anwendung & Rollenspiel: ``Neue Schüler kennenlernen'' & Dialoge mit wechselnden Partnern & PA/GA & Namensschilder \\
\hline

25--35 & Erweiter. & \E{}: Kulturvergleich tú/usted, Variationen & Diskutieren, Notizen & Plenum & Tafel \\
\hline

35--40 & Sicherung & Merkkasten gemeinsam ausfüllen & AB ergänzen & Plenum & AB-01b \\
\hline

40--45 & Abschluss & Zusammenfassung, HA: Vokabeln lernen, LP-01b optional & Notieren & Plenum & -- \\
\hline
\end{longtable}

\textbf{Legende:} EA = Einzelarbeit, PA = Partnerarbeit, GA = Gruppenarbeit

% ============================================================================
% 7. ERWARTETE SCHWIERIGKEITEN
% ============================================================================
\section{Erwartete Schwierigkeiten}

\begin{tabular}{|p{4.5cm}|p{8cm}|}
\hline
\textbf{Schwierigkeit} & \textbf{Reaktion} \\
\hline
Aussprache \textit{eres} ['eres] & Vorsprechen, chorisch wiederholen, Einzelkorrektur \\
\hline
Verwechslung \textit{soy/es} & Kontrastive Übung: ``Ich bin'' vs. ``Er/Sie ist'' \\
\hline
Subjektpronomen vergessen & Hinweis: Im Spanischen optional, aber zur Verdeutlichung hilfreich \\
\hline
Interferenz: *``Yo soy soy'' & Bewusstmachung: Ein Verb reicht! \\
\hline
Zeitproblem & Puffer: Kulturvergleich kürzen, als HA aufgeben \\
\hline
\end{tabular}

% ============================================================================
% 8. HAUSAUFGABE
% ============================================================================
\section{Hausaufgabe}

\begin{itemize}
    \item Lehrwerk S. 12, Übung 3 (ser einsetzen)
    \item Vokabeln lernen: ser-Formen, Redemittel Vorstellung
    \item Optional: Lernpfad LP-01b.html bearbeiten
\end{itemize}

% ============================================================================
% 9. LÖSUNGEN
% ============================================================================
\section{Lösungen}

\subsection{AB-01b Lösungen}

\textbf{Aufgabe 1: Konjugation ser}
\begin{itemize}
    \item yo $\rightarrow$ soy
    \item tú $\rightarrow$ eres
    \item él/ella $\rightarrow$ es
    \item nosotros $\rightarrow$ somos
\end{itemize}

\textbf{Aufgabe 2: Lückentext}
\begin{enumerate}[label=\alph*)]
    \item Me llamo María. \textbf{Soy} de Alemania.
    \item ¿Cómo te llamas? -- \textbf{Soy} Pedro.
    \item Él \textbf{es} de España.
    \item Nosotros \textbf{somos} de Rostock.
\end{enumerate}

\textbf{Aufgabe 3: Dialog (Beispiel)}
\begin{itemize}
    \item A: ¡Hola! ¿Cómo te llamas?
    \item B: Me llamo [Name]. ¿Y tú?
    \item A: Soy [Name]. ¿De dónde eres?
    \item B: Soy de Alemania, de Rostock.
\end{itemize}

\textbf{Aufgabe 4: Freie Produktion (Bewertungskriterien)}
\begin{itemize}
    \item Korrekte ser-Form: 2P
    \item Sinnvolle Redemittel: 2P
    \item Kultureller Bezug (tú/usted): 2P
\end{itemize}

% ============================================================================
% 10. ANLAGEN
% ============================================================================
\section{Anlagen}

\begin{enumerate}
    \item Arbeitsblatt (AB-01b.pdf)
    \item Musterlösung (ML-01b.pdf)
    \item Lehrerhinweise (LH-01b.pdf)
    \item Lernpfad (LP-01b.html)
    \item Redemittel-Karten (Kopiervorlage)
    \item Audio Track 5-6
\end{enumerate}

% ============================================================================
% 11. DIDAKTIK-SCORE
% ============================================================================
\newpage
\section{Didaktik-Score}

\begin{scorebox}
\textbf{Bewertung der didaktischen Qualität nach 10 Dimensionen}\\
\textbf{Ziel-Score: $\geq$ 9.0 (Premium-Qualität)}

\vspace{1em}
\begin{tabular}{|c|l|c|c|p{4cm}|}
\hline
\textbf{\#} & \textbf{Dimension} & \textbf{Gew.} & \textbf{Score} & \textbf{Notiz} \\
\hline
1 & Differenzierung & 15\% & 9/10 & [B]/[S]/[E] qualitativ \\
\hline
2 & Taxonomie (AFB) & 10\% & 9/10 & 25/50/25 erfüllt \\
\hline
3 & Lernziel-Alignment & 10\% & 10/10 & RP MV explizit \\
\hline
4 & Aktivierung & 10\% & 9/10 & Entdeckend, Dialoge \\
\hline
5 & Scaffolding & 10\% & 9/10 & Gestufte Hilfen \\
\hline
6 & Feedback-Qualität & 10\% & 9/10 & LP: elaboriert \\
\hline
7 & Sprachliche Klarheit & 10\% & 10/10 & Kl. 7 angemessen \\
\hline
8 & Motivation & 10\% & 9/10 & Authentisches Audio \\
\hline
9 & Inklusion (UDL) & 10\% & 9/10 & Mehrere Zugänge \\
\hline
10 & Praktikabilität & 5\% & 9/10 & 90 Min. realistisch \\
\hline
\multicolumn{3}{|r|}{\textbf{GESAMT:}} & \textbf{9.2/10} & \\
\hline
\end{tabular}

\vspace{1em}
$\checkmark$ \textcolor{successgreen}{\textbf{FREIGABE} (Score $\geq$ 9.0)}
\end{scorebox}

\end{document}
